\chapter{Restraints}\index{restraints|textbf}\label{Chap:restraints}

While restraints are used for the same purpose as constraints, to fit models that are more complex than what can be managed with the limited set of observations allowed by our powder diffraction measurement(s). However, constraints and restraints work in completely different ways. Constraints remove degrees of freedom from the model. If our model structure has an asymmetric unit\index{asymmetric unit} with five atoms, where none are on special positions\index{special positions} and with no vacancies, then we have 20 structural parameters ($5\times 3$ coordinates plus 5  $\rm U_{iso}$ values). If we create a constraint that groups 4 of the atoms to have the same $\rm U_{iso}$ value, then we replace 4 parameters with 1 and now have only 17 structural parameters, which is a significantly less complex model. The model will also have other non-structural parameters, but they will not directly change reflection intensities, so they are not as directly coupled to the number of observed reflections. 

Restraints function where you provide a target value for something that can be determined from the model, and there is a penalty function that makes the fit worse the further the model refines from that target. Along with the target value for the restraint, you will also provide both an uncertainty and a weight for the restraint values. As was noted in Chapter \ref{Chap:Leastsquares}, the goal of our optimization is to minimize the $\chi^2$ function, 
$$\chi^2 = \sum_j^N w_j(y_{obs,j}-y_{calc,j})^2$$
where $y_{obs,j}$ and $y_{calc,j}$ are the observed and computed powder diffraction intensities, respectively, and $w_j$ are the weights for each intensity, where optimally $w_j = 1/\sigma_j^2$. When we add restraints to this, we add an additional set of terms for the restraints, which can be expressed as:
$$\chi^2 = \sum_j^N w_j(y_{obs,j}-y_{calc,j})^2 + w_{R} \sum_k^M  \left ( \frac{R_k - R_{calc,k}}{\sigma_{R,k}}\right)^2$$
where there are $M$ restraints, $R_k$ are the restraint values, $R_{calc,k}$ is the value we compute for $R_{k}$  from the structural model, $\sigma_{R,k}$ is the uncertainty on the observation and $w_{R}$ is the weight on the restraints. It is worth noting that in practice, other than weighting, the $R_k$ values are handled no differently than diffraction observations -- as far as the math goes. However, in fact restraints are not something we measure. They are values that we believe are true about our material. I will discuss in \S\ref{ref:testrestraints}, below, thoughts on how to test if what we think we know ain't so. 

Why does one use restraints? More complex structural models can have the problem that the fit is just not that sensitive to some of aspects of the structure, so that if the coordinates or occupancies are refined freely, the resulting structural model is implausible in some fashion. A good example for this would be the fitting of zeolites to powder data. 

As a bit of background on zeolites: They are constructed from tetrahedral networks of silicon and aluminum atoms linked by oxygen atoms. Each aluminum will be accompanied by a charge-balancing cation. For some materials, the net scattering from the cations may exceed the net scattering from the framework (Si, Al and O) atoms. Zeolites often have some intrinsic disorder that limits the range of observable reflections, so the range of observed reflections is often limited. It is common that if one tries to fit the coordinates for a zeolite, against even high quality synchrotron or neutron diffraction data, the coordinates will refine to give Si-O and Al-O bond distances that are well outside the expected values and/or O-Si-O/O-Al-O bond angles that are far from tetrahedral. We know well from both plenty of crystal structures, as well as high-level \textit{ab initio} theory that it takes a lot of energy to force Si-O and Al-O bond distances from their preferred values (Si-O 1.62 \AA, Al-O 1.74 \AA) by more than a few hundredths of an \AA. This is not seen in these materials. Bending the bonding angles takes less energy, but I would expect them to be within 5-10$^\circ$ of the tetrahedral angle of 109.5$^\circ$ unless the framework is very strained. 

When we fit a zeolite structure to powder diffraction data, we want not only a structural model that fits the data, we also that our model will have reasonable bond distances and angles. 
In one of my early zeolitic material structures, back when I was an industrial scientist, I measured both high resolution synchrotron and neutron powder diffraction data on Ca LSX (LSX is a synthetic form of the zeolite faujasite, where Al and Si atoms strictly alternate; this is one of the most simple of all zeolite structures).\label{ref:CaLSX}
For reasons that I do not remember, the structure was initially fitted to each set of data separately, before a combined refinement was performed with both together. 
There were serious inaccuracies with each model when the data sets were handled separately. With the x-ray dataset alone, the O atoms did not refine to positions at reasonable distances from the Si and Al atoms. With the neutron dataset there were problems with the Ca atoms. 
The problem with use of the x-ray data alone was not that the fit produced unreasonable Si-O and Al-O distances and angles, it was that there were a continuous range of  structural models that all produced about the same quality of fit to the diffraction data, but many unreasonable bonding geometries. Had only the x-ray data had been available, we would want a model that not only fits the data well, but also has a plausible geometry for the zeolite framework atoms. Use of restraints on the Si-O and Al-O distances, and possibly the O-Si-O and O-Al-O bond angles would have provided such a model. Without the restraints, the details of the framework geometry would be dictated by the starting point for the refinement and perhaps minor statistical scatter in the data.  
For what it is worth, the combined fit to both sets of data was excellent and no restraints were needed. 

In the GSAS-II GUI, access to restraints is found as a child data tree entry, named for each phase, under Restraints. However, if a project has only one phase, clicking on the top level Restraints data tree entry will also open the restraints for that phase. GSAS-II provides seven types of restraints. Note that GSAS-II allows you to provide a separate $w_R$ (weight) value for each type of restraint. The seven types of restraints are discussed individually, below. You select between the different types of restraints by clicking on the notebook tab at the top of the data window. Note that each restraint type has a ``Use'' checkmark button and a ``Restraint weight factor'' entry value. How to select the weight is discussed in \S\ref{ref:weightrestraints}, below. If ``Use'' is not checked, the restraint settings are ignored. If the weight is set to 0, the restraints are evaluated, but do not affect the refinement.

\section{Bond Distance Restraints}
Restraints on bond distances are the most commonly used type of restraints implemented in GSAS-II. The notebook tab for them is labeled ``Bond'' and they operate on atomic coordinates. They are most commonly used for keeping distances between bonded atoms at a target distance, but they can be used on any pair of atoms in a structure, whether bonded or not. 

To create a bond distance constraint, you must search for interatomic distances. This is done in a series of steps. Before starting you should have an idea of the bonding environment for each type atom that you wish to constrain and what the current distances range between those atoms. If needed, go to the Phase and use the Atoms tab and the Compute/``Show Distances \& Angles'' menu command. You may need to change the radius values for atoms. Note that the results are listed on the console window. 

Then, select the ``Bond'' tab, and make note of the ``Search range'' factor. This will be used to find distances. You will supply a target distance and distances that are up to the length of the Search range'' times the ``target distance'' will be considered as potential restraints. Pairs of atoms with distances longer than that will be ignored. 
You will then use the  ``Edit Restr.''/``Add restraint'' command and then you will be asked to 
\begin{itemize}
    \item in the first window you will supply the ``origin'' atom(s) in the first window. 
    \item in the second window, you will be asked for the ``target'' atoms. These can be the same as the origin atoms. 
    \item in the third window you will supply the target distance.
\end{itemize}
Once this is done, distance constraints are placed into the table of constraints shown in the data window. You can edit the target and esd ($\sigma$) value directly in the table. Or you can select restraints in the table and change the target or esd values for the selected restraints using ``Edit Restr.'' menu commands. There is also a menu command that will delete the selected restraints. If no selection has been made in the table before entering any of these the menu commands, then a window will be displayed where restraints can be selected. 

As an example for how the origin and target atoms work in searches, consider a structure with Si, O and Al atoms. If the origin is set as the O atoms and the target is set to the Si atoms, then only distances between O and Si atoms will be considered. If both O are Si are selected for both the target and the origin, then O-O, Si-Si, O-Si and Si-O distances will be considered, where the latter two distances are duplicates. There is no problem with including the same pair of atoms twice, but that distance will have double weight compared to one that is only listed once. 

\subsection{Import from Mogul}
The Cambridge Crystallographic Data Center produces a program called Mercury\index{software!Mercury} that tabulates distances and angles within a molecule as well as prepares the statistically expected values for those distances from the structures in the Cambridge Crystallographic Database using the CSD-Core/Mogul Geometry Check\index{software!MOGUL}. This is primarily of use for organic moieties in structures. See the GSAS-II Help page on Restraints for more information on this command. When used from the Distance tab, constraints will be created for  distances from this file. 

\section{Angle Restraints}
Angle restraints also operate on atomic coordinates. The notebook tab for them is labeled ``Angle''. 
They are most commonly used for keeping angles between bonded atoms at a target, but they can be used on any pair of atoms in a structure, whether bonded or not. As was noted for bond distance constraints, you should have an idea of the bonding environment for each type atom that you wish to constrain. Since the searching is done using atomic radii, it is worthwhile to the Phase entry in the data tree and use the Atoms tab and the Compute/``Show Distances \& Angles'' menu command and confirm that the radii settings locate the angles you wish to constraint. You may need to change the radius values for atoms. Note that the results are listed on the console window. Note that the atoms in an angle are the outer or A/C atom (where the angle is listed as A-B-C) and the central or B atom. 

Then, select the ``Angle'' tab, and make note of the ``Search range'' factor. This will be used to find angles for atoms where the A-B and C-B distances are shorter than the sum of the atom radii distances times the Search range''. Pairs of atoms with distances longer than that will be ignored. 
You will then use the  ``Edit Restr.''/``Add restraint'' command and then you will be asked to: 
\begin{itemize}
    \item in the first window name you will supply the outer atom(s) in the angle. 
    \item in the second window, you will be asked for the central atoms. These can be the same as the origin atoms. 
    \item in the third window you will provide the target angle.
\end{itemize}
Once this is done, located angles are used to create constraints that are placed into the table of constraints shown in the data window. You can edit the target and esd ($\sigma$) value directly in the table. Or you can select restraints in the table and change the target or esd values for the selected restraints using ``Edit Restr.'' menu commands. There is also a menu command that will delete the selected restraints. If no selection has been made in the table before entering any of these the menu commands, then a window will be displayed where restraints can be selected. 

\subsection{Import from Mogul}
As noted, the Mercury program that tabulates distances and angles within a molecule also prepares the statistically expected values for those values from the structures in the Cambridge Crystallographic Database using the CSD-Core/Mogul Geometry Check. This is primarily of use for organic moieties in structures. See the GSAS-II Help page on Restraints for more information on this command. When used from the Angles tab, constraints will be created for  angles from this file. 

\section{Plane Restraints}
Plane restraints are used to restrain atoms to lie on a plane. The notebook tab for them is labeled ``Plane''. One selects four or more atoms and the least-squares plane is computed for those atoms and the root-mean-square distance of the atoms from that plane. This is mostly used in macromolecular crystallography. I have never needed this constraint. 

\section{Chiral Restraints}
Chiral restraints create a chiaral volume constraint. I don't think this has any use outside macromolecular crystallography. 

\section{Chemical Composition Restraints}
This is used to place restraints on occupancies and is found with the notebook tab labeled ``Chem. comp.'' This probably the next most commonly used restraint after bonds and is quite versatile. This constraint uses the number of atoms in the unit cell, which is determined from the site multiplicity times the occupancy. 

The first step in creating a composition constraint is to select the ``Chem. comp.''  tab and then use the  ``Edit Restr.''/``Add restraint'' menu command. The first window to be opened allows selection of which atoms to include into the constraint. What you will select here will depend on what you want the constraint to do. In the next window you specify the target value that will be used in the constraint. (This is easily changed later). The constraint is then created as a list of the selected atoms. For each atom, there are four entries: the atom name, the number of atoms in the cell (the multiplicity * occupancy), a factor that you supply (which defaults as 1) and the product of the factor times the number of atoms. The sum of these products is the calculated value for the restraint, and this is shown as the last line for the restraint, along with the target and $\sigma$ (esd) value. Note that the ``Edit Restr.''/``Delete restraint'' menu command can be used to either delete an entire restraint or remove one or more atom from a restraint, depending on which row(s) are selected. Rows can be selected before the menu command is used by clicking on row labels or if no entries are selected, a window will be opened to select row(s). 

How this restraint will work will depend on what atoms are included and how the factor and target values are set. Here are some ideas: 
\begin{itemize}
    \item Restrain the number of atoms of a type in the unit cell: choose all of those atoms, set the factor to 1 and the target to the specified number of atoms.
    \item Build a charge balance restraint for site substitution, where you include all of the atoms for the types that are substituting for each other, the oxidation state for each atom and the total number of charges expected.
    \item Construct a valence balancing restraint by including all atoms and setting the factor as each atom's valence (positive and negative). The next charge should be zero. 
    \item Restrain the mass of the unit cell to a target by setting the factors for each atom to the atomic weight. 
\end{itemize}
You can create as many compositional restraints as you wish. 

\section{Magnetic Moment Restraints}
The notebook tab for them is labeled ``Moments'' and this tab will only appear for magnetic phases. They operate only on the magnetic moments for atoms. This restraint requires inclusion of two or more magnetic atoms and restraints the included atoms to all have the same moment magnitude, independent of spin direction. This is what would be expected, as the spin magnitude is dictated by the atom type and should not vary due to the site. The target value is the mean moment magnitude and cannot be edited. 

\section{General Restraints}
I created this restraint because I felt it might someday be needed, but I have yet to find a problem where it has been deployed. This option allows you to write your own equation that will be evaluated after substituting in values for GSAS-II parameters. 

To create a general restraint, one selects the ``General'' tab and then uses the ``Edit Restr.''/``Add restraint'' menu command or the ``Add restraint'' button on the window. That opens a window where an expression can be entered as a single line of Python. That Python command can invoke math, numpy or GSAS-II routines or even your own external library of Python code. If that Python expression contains any variables, you will then assign GSAS-II parameters by name (see \S\ref{ref:paramname}) to each of those variable. As an example, if the expression 

\begin{quote}
\texttt{10*x**2 + b*x + c}    
\end{quote}
is entered, then you will be asked to supply GSAS-II parameters of form \texttt{p:h:name} for \texttt{x}, \texttt{b} and \texttt{c}. When evaluating $\chi^2$, the value for those parameters will be substituted into the expression and the result is placed as $R_{calc,k}$. After creating the General restraint, the target value for the restraint, the $\sigma$ value (labeled as esd) and the weight factor are entered in the data window. 

\section{Choosing the Restraint Weighting}\label{ref:weightrestraints}
The weight factor to be used in the fit for restraints is not a simple choice and the value you will want may well change over the course of a fit. It should be noted in the sum for $\chi^2$,
$$\chi^2 = \sum_j^N w_j(y_{obs,j}-y_{calc,j})^2 + w_{R} \sum_k^M  \left ( \frac{R_k - R_{calc,k}}{\sigma_{R,k}}\right)^2$$
there are $N$ measured powder diffraction data points and $M$ restraints. $N$ is usually several orders of magnitude larger than $M$, so if $w_R$ is left at the default value, the diffraction data will usually contribute a vastly larger share of the total $\chi^2$ than the restraints. This is usually the opposite of what I want at the initial stages of a refinement. For this reason, when I first start refining structural parameters, I want to prevent the model from distorting into some unreasonable way, so it is common for me to initially set the restraints weight to a very large number, on the order of $N$ or even $10N$ ($10^3$ to  $10^5$). One can estimate the impact of the weighted constraints by looking at the amount that each term in the sum contributes to the total $\chi^2$ from values that are displayed on the console as the refinement progresses; these are also shown in the GUI in some places. 

One the fit to the data has progressed, I will start lowering the $w_R$ value(s) and watch what happens to the quantities that are being constrained (bond distances, etc.). Ideally, I can lower $w_R$ to 0 and still have only small deviations from the ideal distances, etc. This is ideal, because that means that the restraint was only needed to guide the refinement to a structurally realistic minimum, but that local minimum (and one hopes that is the global minimum) is stable. I tend to leave the restraints in place, even with zero weight so that I can keep an eye on the fitted values and confirm that they stay reasonable. However, if a zero weight is not possible, then frequently a value of $w_R$ in the range of 1 to 10 is all that is needed to push the model to have a small restraint contribution that keeps the values in the expected ranges. Do make a note of the restraint contribution to the total $\chi^2$ and report that when you describe how the model was fit. When publishing, do not report the restrained quantities with their least-squares uncertainties as if they were determined by the fit. Those quantities have been biased by the restraints and this needs to be clear to your readers. 

\section{Testing Restraints}\label{ref:testrestraints}
Since it is possible to code a restraint that is an incorrect assumption, we need to understand what the impact of that will be on our fit, with the hope that it can be detected. The key concept here is that if a restraint is realistic, then its inclusion in the fit will change the fit to the powder diffraction data only by very minor amounts (this is best noted by the $\rm R_{wp}$ value for the powder pattern). It may have a large impact on the values of the restrained parameters, but as the weighting is reduced for the restraint, the $\rm R_{wp}$ value should only change by perhaps 0.1 or less.  

To understand this, as an example, let us consider a material composed of tetrahedral Al and Si atoms linked by O atoms. As was noted before, the expected bond distances for Al and Si are a bit different, 1.74 \AA\  and 1.62 \AA, respectively. Suppose that I make a mistake and label a site that actually contains an Al atom as having a Si atom and place restraints on the bond distances around that atom as 1.62 \AA, even though the structure really has the longer Al-O distances. This incorrect model will have two inaccuracies. One is that there is an extra electron being placed on that site, but experience tells me that this will have negligible impact on the fit. The $\rm U_{iso}$ will be enlarged a bit, but the $\rm R_{wp}$/GOF values will stay about the same. However, if the data have sufficient sensitivity to actually determine the bond distances for that site, the data and restraints will be in conflict. Without restraints, the bond distances would average about 1.74 \AA, but restraints will push the values to be shorter. As the weighting of the restraints are lowered, the $\rm R_{wp}$ value should drop as the fit moves away from the mistaken distances and towards the correct values. This may be be a subtle difference if the impact of the bond distances is very small. Also, if the restraint is needed because the model is insensitive to the bond distances, then one might expect the range of distances to become larger, but the average should stay close to the target, but if the target is wrong, then the average will change as the weighting is lowered. In truth, this mistake might be too subtle to find, but knowing what to look for at least makes it possible. With more egregious mistakes, the conflict between the data and the restraints would be more obvious. 

\section{Restraints not provided by GSAS-II}
GSAS-II does not provide a restraint to keep $\rm U_{iso}$ values positive or for that matter from refining to values that are too large (which is worse). This is not an oversight. Implausible $\rm U_{iso}$ values are telling you something about your data and/or your model that you do not want hidden by a restraint. Either the data are not sensitive to the $\rm U_{iso}$ values that refine to unbelievable numbers. This should not be hidden by a restraint. In that case the values should be set and not refined, or grouped so that the value has more leverage. Alternatively, the data are indeed sensitive to those $\rm U_{iso}$ values, in which case your data is sending you a subtle message that there is something wrong with your model. Listen to that. 
